\section{SD-C14 and its determinant data types}
\label{sec:determinants}

Let $F_p$ be the tensor-prime full-shift atom with entropy $\log p$.  Its
primitive loop $\gamma_p$ and repetitions $\gamma_p^r$ remain the base
grammar.  The fiber is
\[
 \cA=\CC\oplus L(\ZZ),\qquad W=1\oplus u,                 \tag{4.1}
\]
where the canonical group trace satisfies $\tau(u^r)=0$ for $r\ne0$.  For a
frozen $c\ge0$, put
\[
 \Phi_c(a\oplus x)=a+c\tau(x).                             \tag{4.2}
\]
For $c>0$ this is a faithful finite positive trace, but
\[
 \Phi_c(1)=1+c,\qquad \Phi_c(W^r)=1\quad(r\ne0).
 \tag{4.3}\label{eq:power-moments}
\]
The normalized trace $\widetilde\Phi_c=\Phi_c/(1+c)$ instead gives
$\widetilde\Phi_c(W^r)=1/(1+c)$, so normalization and the exact one-copy
ledger are incompatible when $c>0$.

\subsection{The analytic trace-log}

For $|q|<1$, norm functional calculus defines
\[
 D_c(q)=\exp\Phi_c\bigl(\log(1-qW)\bigr).                 \tag{4.4}
\]

\begin{theorem}[Exact ghost blindness]
\label{thm:blindness}
For every $c\ge0$ and $|q|<1$,
\[
 D_c(q)=1-q.                                               \tag{4.5}
\]
More generally, two scalar trace-log germs agree near $q=0$ if and only if
their aggregated positive power traces agree at every order.
\end{theorem}

\begin{proof}
The logarithm series and the moment identity in
\eqref{eq:power-moments} give
\[
 \log D_c(q)
 =-\sum_{r\ge1}\frac{q^r}{r}\Phi_c(W^r)
 =-\sum_{r\ge1}\frac{q^r}{r}=\log(1-q).                  \tag{4.6}
\]
Equality of germs is equivalent to equality of their Taylor coefficients,
which are the power traces divided by $r$.
\end{proof}

On the tensor-prime direct sum, take $q=zp^{-s}$ at atom $p$.  For
$\RePart s>1$ and $|z|\,2^{-\RePart s}<1$, absolute convergence yields
\[
 D_c(s,z)=\prod_pD_c(zp^{-s})=\prod_p(1-zp^{-s}),         \tag{4.7}
\]
independently of $c$.  This is an exact analytic determinant and an unchanged
Euler ledger.  It contains no new Haar-fiber information.

\subsection{Magnitude and normalized spectral distribution}

The Fuglede--Kadison determinant answers a different question.

\begin{proposition}[Exact magnitude formula]
\label{prop:fk}
For $q\in\CC$,
\[
 \Delta_{\Phi_c}(1-qW)
 =|1-q|\max(1,|q|)^c.                                     \tag{4.8}
\]
\end{proposition}

\begin{proof}
Spectral calculus and Jensen's circle formula give
\begin{align}
 \log\Delta_{\Phi_c}(1-qW)
 &=\log|1-q|+c\int_\TT\log|1-qz|\,dm_{\rm H}(z)\\
 &=\log|1-q|+c\log^+|q|.                                 \tag{4.9}
\end{align}
Exponentiation proves the formula.
\end{proof}

Inside $|q|<1$, \cref{prop:fk} is as blind as the analytic determinant.
Outside the disk it detects $c$ through radial magnitude, but is positive
real rather than a holomorphic oriented divisor.  With the normalized trace,
the normal unitary $W$ has probability spectral distribution
\[
 \frac1{1+c}\delta_1+\frac{c}{1+c}m_{\rm H}.              \tag{4.10}
\]
It sees the continuous circle spectrum only after diluting the base atom
coefficient.  These coordinates cannot be spliced into \cref{thm:blindness}
as a single Route-A determinant.
